\documentclass{article}

\usepackage[OT2, T1, TS1, X2, ECMSRB1, ECMSRB2, T2A]{fontenc}
\usepackage[utf8]{inputenc}
\usepackage[serbianc]{babel}
\usepackage[a4paper,top=2.0cm,left=3cm,right=2.5cm,bottom=2cm,includefoot,includehead]{geometry}

\usepackage{cmsrb}
\usepackage{cmupint}
\let\enint\intop

\usepackage{amssymb}

\input glyphtounicode.tex
\pdfgentounicode=1

\usepackage{fancyvrb}
\usepackage{tabularx}

\newcommand\textcmsuper[1]{{\fontencoding{T2A}\fontfamily{cmr}\selectfont #1}}
\newcommand{\example}[1]{{\fontencoding{T2A}\selectfont #1} & \textcmsuper{#1}}

\newcommand\otinput[1]{{\fontencoding{T1}\selectfont \verb"#1"} & {\fontencoding{OT2}\selectfont #1}}
\newcommand\ecminput[1]{{\fontencoding{T1}\selectfont \verb"#1"} & {\fontencoding{ECMSRB1}\selectfont #1} & {\fontencoding{ECMSRB2}\selectfont #1}}

\DeclareTextSymbolDefault{\dj}{T1}

\newcounter{primer}
\newcounter{rezultat}
\newenvironment{primer}{\VerbatimEnvironment\refstepcounter{primer}\par\medskip
	\noindent\textbf{Пример~\theprimer: }
\begin{Verbatim}
}{\end{Verbatim}\medskip}
\newenvironment{rezultat}{\refstepcounter{rezultat}\par\medskip
\noindent\textbf{Резултат~\therezultat: }\\ }{\medskip}

\makeatletter
\newcommand{\manuallabel}[2]{\def\@currentlabel{#2}\label{#1}}
\makeatother

\usepackage{fonttable}

\begin{document}
	\title{Пакет \texttt{cmsrb}}
	\author{Урош Стефановић\footnote{\texttt{urostajms@gmail.com}}}
	\date{\today{} верзија 4.0}
	\maketitle
	
	\section{Зашто \textsf{cmsrb}?}
	
	Пакет \textsf{cm-super} обезбеђује одличну подршку за ћирилицу у разним језицима, али код њега постоји проблем са курзивним облицима (италиком) неких слова за српски и македонски језик. Пакет \textsf{cmsrb} садржи исправан облик курзивних слова \verb|г|, \verb|д|, \verb|п|, \verb|т|, и слова \verb|б|. Такође, овај пакет доноси одређена побољшања у словима и акцентима који се користе у српском језику.
	
	\section{Карактеристике пакета}
	
	Фонтови \textsf{cmsrb} су проширење \textit{Computer Modern} фонтова за српски и македонски језик за \TeX{} (\LaTeX). Подржана кодирања су \texttt{T1}, \texttt{TS1}, \texttt{T2A}, \texttt{X2} и \texttt{OT2}, као и експериментална кодирања \texttt{ECMSRB1} и \texttt{ECMSRB2}. Подржани стилови фонта су серифни, санс-серифни и непропорционални, за величину 10\textit{pt}.
	
	Овај пакет је једноставан за коришћење: само је потребно да ставимо
	\begin{verbatim}
	\usepackage{cmsrb}
	\end{verbatim}
	у преамбулу документа.
	
	\begin{table}[h!]
		\newcolumntype{C}{>{\centering\arraybackslash}X}%
		\begin{tabularx}{\textwidth}{|C C|C C|C C|}
			\hline
			cmsrb & cm-super & cmsrb & cm-super & cmsrb & cm-super  \\ \hline
			\hline
			\example{\textit{бгдпт}} & \example{\textbf{\textit{бгдпт}}} & \example{ђћ} \\ \hline
			\example{б\textbf{б}\textit{б}\textsl{б}} & \example{\={а}\textsc{\={а}\f{и}}} & \example{\textit{đ\textbf{đ}}} \\ \hline
		\end{tabularx}
		\caption{Неке \textsf{cm-super} и \textsf{cmsrb} разлике.}\label{t2}
	\end{table}
	
	\section{Карактеристике \textit{OT2} кодирања}
	
	Кодирање \texttt{OT2} је измењено како би подржало правилно пребацивање из латинице у ћирилицу у српском језику. Према томе, \textsf{ts}, \textsf{kh}, \textsf{ch} и сличне лигатуре су уклоњене из кодирања (види табелу~\ref{t1}). 
	
	У \texttt{OT2} кодирању су додата македонска слова Ѓ и Ќ, црногорска слова \CYRSJE{} и \CYRZJE, као и одређени правописни знаци.
	
	Ипак, није препоручљиво коришћење \texttt{OT2} кодирања; бољи избор за српски језик је \texttt{T2A} кодирање и \textit{utf8} унос. Кодирање \texttt{OT2} је добар избор ако већ имамо документ написан латиницом — тада ће бити једноставно пребацити га у ћирилицу.
	
	\begin{table}
	\newcolumntype{C}{>{\centering\arraybackslash}X}%
	\begin{tabularx}{\textwidth}{|C C|C C|C C|C C|}
	\hline
	Улаз & Излаз & Улаз & Излаз & Улаз & Излаз & Улаз & Излаз \\ \hline
	\hline
	\otinput{A} & \otinput{B} & \otinput{C} & \otinput{D} \\ \hline
	\otinput{E} & \otinput{F} & \otinput{G} & \otinput{H} \\ \hline
	\otinput{I} & \otinput{J} & \otinput{K} & \otinput{L} \\ \hline
	\otinput{M} & \otinput{N} & \otinput{O} & \otinput{P} \\ \hline
	\otinput{Q} & \otinput{R} & \otinput{S} & \otinput{T} \\ \hline
	\otinput{U} & \otinput{V} & \otinput{W} & \otinput{X} \\ \hline
	\otinput{Y} & \otinput{Z} & \otinput{\#} & \otinput{} \\ \hline
	\hline
	\otinput{a} & \otinput{b} & \otinput{c} & \otinput{d} \\ \hline
	\otinput{e} & \otinput{f} & \otinput{g} & \otinput{h} \\ \hline
	\otinput{i} & \otinput{j} & \otinput{k} & \otinput{l} \\ \hline
	\otinput{m} & \otinput{n} & \otinput{o} & \otinput{p} \\ \hline
	\otinput{q} & \otinput{r} & \otinput{s} & \otinput{t} \\ \hline
	\otinput{u} & \otinput{v} & \otinput{w} & \otinput{x} \\ \hline
	\otinput{y} & \otinput{z} & \otinput{+} & \otinput{} \\ \hline
	\hline
	\otinput{C1} & \otinput{D1} & \otinput{D2} & \otinput{D3} \\ \hline
	\otinput{E0} & \otinput{E1} & \otinput{E2} & \otinput{I0} \\ \hline
	\otinput{I1} & \otinput{J1} & \otinput{J2} & \otinput{L1} \\ \hline
	\otinput{N0} & \otinput{N1} & \otinput{P1} & \otinput{P2} \\ \hline
	\otinput{Z1} & \otinput{\v{C}} & \otinput{\'C} & \otinput{\DJ} \\ \hline
	\otinput{\v{S}} & \otinput{\v{Z}} & \otinput{LJ} & \otinput{Lj} \\ \hline
	\otinput{NJ} & \otinput{Nj} & \otinput{D\v{Z}} & \otinput{D\v{z}} \\ \hline
	\hline
	\otinput{c1} & \otinput{d1} & \otinput{d2} & \otinput{d3} \\ \hline
	\otinput{e0} & \otinput{e1} & \otinput{e2} & \otinput{i0} \\ \hline
	\otinput{i1} & \otinput{j1} & \otinput{j2} & \otinput{l1} \\ \hline
	\otinput{\i} & \otinput{n1} & \otinput{p1} & \otinput{p2} \\ \hline
	\otinput{z1} & \otinput{\v{c}} & \otinput{\'c} & \otinput{\dj} \\ \hline
	\otinput{\v{s}} & \otinput{\v{z}} & \otinput{lj} & \otinput{} \\ \hline
	\otinput{nj} & \otinput{} & \otinput{d\v{z}} & \otinput{} \\ \hline
	\hline
	\otinput{\char20} & \otinput{\char21} & \otinput{\char28} & \otinput{\char29} \\ \hline
	\otinput{\'G} & \otinput{\'K} & \otinput{\'g} & \otinput{\'k} \\ \hline
	\otinput{\'S} & \otinput{\'Z} & \otinput{\'s} & \otinput{\'z} \\ \hline
	\end{tabularx}
	\caption{Кодирање \texttt{OT2} за пакет \textsf{cmsrb}.}\label{t1}
	\end{table}

\begin{primer}
\documentclass{article}
\usepackage{cmsrb}
\usepackage[OT2,T1]{fontenc}
\usepackage[serbian]{babel}
\newcommand{\test}%
{Ljubazni fenjerd\v zija \v ca\dj avog lica ho\'ce da mi poka\v ze \v stos.}
\begin{document}
\test \\
\fontencoding{OT2}\selectfont \test \\
Akcenti: \'a\`a\C a\f a\=a\^a\"a\u a
\end{document}
\end{primer}

\begin{rezultat}
\indent {\fontencoding{T1}\selectfont Ljubazni fenjerd\v zija \v ca\dj avog lica ho\'ce da mi poka\v ze \v stos.} \\
{\fontencoding{OT2}\selectfont Ljubazni fenjerd\v zija \v ca\dj avog lica ho\'ce da mi poka\v ze \v stos. \\ Akcenti: \'a\`a\C a\f a\=a\^a\"a\u a}
\end{rezultat}

\begin{primer}
\documentclass{article}
\usepackage{cmsrb}
\usepackage[OT2,T1]{fontenc}
\usepackage[utf8]{inputenc}
\usepackage[serbian]{babel}
\newcommand{\test}%
{Ljubazni fenjerdžija čađavog lica hoće da mi pokaže štos.}
\begin{document}
\textit{\test} \\
\fontencoding{OT2}\selectfont \textit{\test}
\end{document}
\end{primer}

\begin{rezultat}
\indent{\fontencoding{T1}\selectfont \textit{Ljubazni fenjerd\v zija \v ca\dj avog lica ho\'ce da mi poka\v ze \v stos.}} \\
{\fontencoding{OT2}\selectfont \textit{Ljubazni fenjerd\v zija \v ca\dj avog lica ho\'ce da mi poka\v ze \v stos.}}
\end{rezultat}

	\section{Карактеристике \textit{T1} кодирања}
	
	Кодирање \texttt{T1} сада подржава конверзију из ћирилице у латиницу (види пример~\ref{prim:t1}).
	
	На позицији 208 се налази слово {\fontencoding{T1}\selectfont Đ (\texttt{U+0110})} уместо слова {\fontencoding{TS1}\selectfont \DH{}  \fontencoding{T1}\selectfont (\texttt{U+00D0})}, чиме се омогућава копирање или претрага текста написаног на латиници.
	
\begin{primer}
\documentclass{article}
\usepackage{cmsrb}
\usepackage[T2A,T1]{fontenc}
\usepackage[utf8]{inputenc}
\usepackage[serbian]{babel}
\begin{document}
Љубазни фењерџија чађавог лица хоће да ми покаже штос.
\end{document}
\end{primer}
\manuallabel{prim:t1}{\theprimer}

\begin{rezultat}
\indent Ljubazni fenjerdžija čađavog lica hoće da mi pokaže štos.
\end{rezultat}
	
	
	\section{Карактеристике \textit{T2A} кодирања}
	
	Ћириличка слова \texttt{Ј} и \texttt{ј} нису подржана у \texttt{T2A} кодирању; због тога су она постављена уместо латиничких слова \texttt{J} и \texttt{j}.
	То даје бољи резултат у ћириличком тексту, али то такође значи да се та латиничка слова не налазе у овом кодирању, па је потребно водити рачуна приликом копирања и претраге латиничког текста написаног \texttt{T2A} кодирањем.
	
	Слова {\fontencoding{X2}\selectfont \char"81, \char"89, \char"86, \char"91, \char"A1, \char"A9, \char"A6, \char"B1} су уклоњена из кодирања \texttt{T2A}, а уместо њих су стављена слова \char"81, \char"89, \char"86, \char"91, \char"A1, \char"A9, \char"A6, \char"B1.

\begin{primer}
\documentclass{article}
\usepackage{cmsrb}
\usepackage[T2A]{fontenc}
\usepackage[utf8]{inputenc}
\usepackage[serbianc]{babel}
\DeclareTextSymbolDefault{\dj}{T1}
\begin{document}
\textit{Ljubazni fenjerdžija čađavog lica hoće da mi pokaže štos.\\
Љубазни фењерџија чађавог лица хоће да ми покаже штос.} \\
Акценти: \'{и}\`{и}\C{и}\f{и}\={и}\^{и}\"{и}\u{и}
\end{document}
\end{primer}

\begin{rezultat}
\indent {\textit{Ljubazni fenjerdžija čađavog lica hoće da mi pokaže štos.\\ Љубазни фењерџија чађавог лица хоће да ми покаже штос.} \\ Акценти: \'{и}\`{и}\C{и}\f{и}\={и}\^{и}\"{и}\u{и}}
\end{rezultat}

\begin{primer}
\documentclass{article}
\usepackage{cmsrb}
\usepackage[T2A]{fontenc}
\usepackage[utf8]{inputenc}
\usepackage[serbianc]{babel}
\begin{document}
\textit{абвгдђѓежз\'{з}ѕијклљмнњопрс\'{с}тћќуфхцчџш}
\end{document}
\end{primer}

\begin{rezultat}
\indent {\textit{абвгдђѓежз\'{з}ѕијклљмнњопрс\'{с}тћќуфхцчџш}}
\end{rezultat}
	
	\section{Карактеристике \textit{TS1} кодирања}
	
	У \texttt{TS1} кодирању се налази слово {\fontencoding{TS1}\selectfont \DH} (\texttt{U+00D0}) које је уклоњено из \texttt{T1} кодирања.
	Такође, у \texttt{TS1} кодирању се налазе слова б, г, д, п, т са руским обликом курзива.

\begin{primer}
\documentclass{article}
\usepackage{cmsrb}
\usepackage[TS1,T2A]{fontenc}
\usepackage[utf8]{inputenc}
\usepackage[serbianc]{babel}
\begin{document}
\textit{бгдпт} \fontencoding{TS1}\selectfont \textit{бгдпт \DH}
\end{document}
\end{primer}

\begin{rezultat}
	\indent \textit{бгдпт} {\fontencoding{TS1}\selectfont \textit{бгдпт \DH}}
\end{rezultat}
	
	\section{Математика}
	
	Пакет \textsf{cmsrb} учитава подразумеване \textit{Computer Modern} математичке фонтове, али учита само величине веће или једнаке 10\textit{pt}, за бољи визуелни ефекат.
	Ако се користи \textsf{math} опција тада ће бити учитан пакет \textsf{cmupint} који мења подразумевани знак за интеграл $\enint$ у $\int$ (усправни интеграл је традиционално коришћен у српском језику, са \verb|\limits| опцијом).
	Такође, у српском језику боље је коришћење знакова $\leqslant$ и $\geqslant$ уместо знакова $\leq$ и $\geq$. Ако се користи \textsf{math} опција учитаће се пакет \textsf{amssymb} и ти знаци ће аутоматски бити промењени (исто важи за $\nleqslant$ и $\ngeqslant$).
	
	Опција се користи на следећи начин:
	\begin{verbatim}
	\usepackage[math]{cmsrb}
	\end{verbatim}
	
\begin{primer}
\documentclass{article}
\usepackage{cmsrb}
\newcommand{\ud}{\,\mathrm{d}}
\begin{document}
$$ \int_0^1 e^x\ud x \geq 0 $$
\end{document}
\end{primer}

\begin{rezultat}
$$ \enint\nolimits_0^1 e^x\,\mathrm{d} x \geq 0 $$
\end{rezultat}

\begin{primer}
\documentclass{article}
\usepackage[math]{cmsrb}
\newcommand{\ud}{\,\mathrm{d}}
\begin{document}
$$ \int_0^1 e^x\ud x \geq 0 $$
\end{document}
\end{primer}

\begin{rezultat}
$$ \int_0^1 e^x\,\mathrm{d} x \geqslant 0 $$
\end{rezultat}

	
	\section{Експериментална кодирања}
	
	Кодирања \texttt{ECMSRB1} и \texttt{ECMSRB2} су направљена само за овај пакет. Слична су са \texttt{OT2} кодирањем, али садрже специфичне лигатуре и акценте (види табелу~\ref{t3}). За комплетну табелу карактера који су подржани овим кодирањем погледати табеле на странама~\pageref{ecs1} и \pageref{ecs2}. Основна намена ових кодирања је писање текста на српском језику само \texttt{ASCII} карактерима, као и једноставно пребацивање из латинице у ћирилицу.
	
	\begin{table}[h!]
		\newcolumntype{C}{>{\centering\arraybackslash}X}%
		\begin{tabularx}{\textwidth}{|C C C|C C C|C C C|}
			\hline
			Улаз & \texttt{ECMSRB1} & \texttt{ECMSRB2} & Улаз & \texttt{ECMSRB1} & \texttt{ECMSRB2} &Улаз & \texttt{ECMSRB1} & \texttt{ECMSRB2} \\ \hline
			\hline
			\ecminput{CX} & \ecminput{CY} & \ecminput{DX} \\ \hline
			\ecminput{DY} & \ecminput{EY} & \ecminput{GY} \\ \hline
			\ecminput{KY} & \ecminput{LJ} & \ecminput{LY} \\ \hline
			\ecminput{NJ} & \ecminput{NY} & \ecminput{SX} \\ \hline
			\ecminput{SXY} & \ecminput{ZX} & \ecminput{ZY} \\ \hline
			\ecminput{ZXY} & \ecminput{\EZH} & \ecminput{} \\ \hline
			\hline
			\ecminput{Cx} & \ecminput{Cy} & \ecminput{Dx} \\ \hline
			\ecminput{Dy} & \ecminput{Ey} & \ecminput{Gy} \\ \hline
			\ecminput{Ky} & \ecminput{Lj} & \ecminput{Ly} \\ \hline
			\ecminput{Nj} & \ecminput{Ny} & \ecminput{Sx} \\ \hline
			\ecminput{Sxy} & \ecminput{Zx} & \ecminput{Zy} \\ \hline
			\ecminput{Zxy} & \ecminput{} & \ecminput{} \\ \hline
			\hline
			\ecminput{cx} & \ecminput{cy} & \ecminput{dx} \\ \hline
			\ecminput{dy} & \ecminput{ey} & \ecminput{gy} \\ \hline
			\ecminput{ky} & \ecminput{lj} & \ecminput{ly} \\ \hline
			\ecminput{nj} & \ecminput{ny} & \ecminput{sx} \\ \hline
			\ecminput{sxy} & \ecminput{zx} & \ecminput{zy} \\ \hline
			\ecminput{zxy} & \ecminput{\ezh} & \ecminput{} \\ \hline
			\hline
			\ecminput{\`a} & \ecminput{\'a} & \ecminput{\^a} \\ \hline
			\ecminput{\T a} & \ecminput{\D a} & \ecminput{\H a} \\ \hline
			\ecminput{\r a} & \ecminput{\R a} & \ecminput{\v a} \\ \hline
			\ecminput{\u a} & \ecminput{\U a} & \ecminput{\=a} \\ \hline
			\ecminput{\b a} & \ecminput{\.a} & \ecminput{\d a} \\ \hline
			\ecminput{\c a} & \ecminput{\k a} & \ecminput{\"a} \\ \hline
			\ecminput{\C a} & \ecminput{\~a} & \ecminput{\f a} \\ \hline
		\end{tabularx}
		\caption{Кодирања \texttt{ECMSRB1} и \texttt{ECMSRB2} за пакет \textsf{cmsrb}.}\label{t3}
	\end{table}

\begin{primer}
\documentclass{article}
\usepackage{cmsrb}
\usepackage[ECMSRB1]{fontenc}
\usepackage[serbian]{babel}
\begin{document}
Ljubazni fenjerdxija cxadyavog lica hocye da mi pokazxe sxtos.
\end{document}
\end{primer}

\begin{rezultat}
	\indent {\fontencoding{ECMSRB1}\selectfont Ljubazni fenjerdxija cxadyavog lica hocye da mi pokazxe sxtos.}
\end{rezultat}

\begin{primer}
\documentclass{article}
\usepackage{cmsrb}
\usepackage[ECMSRB2]{fontenc}
\usepackage[serbian]{babel}
\begin{document}
Ljubazni fenjerdxija cxadyavog lica hocye da mi pokazxe sxtos.
\end{document}
\end{primer}

\begin{rezultat}
\indent {\fontencoding{ECMSRB2}\selectfont Ljubazni fenjerdxija cxadyavog lica hocye da mi pokazxe sxtos.}
\end{rezultat}
    
    \section{Историја верзија}
    
    \subsection*{4.0}
    
    \begin{itemize}
    	\item Све опције уклоњене, опција \textsf{math} додата.
    	\item Измењена кодирања \texttt{T1}, \texttt{T2A}, \texttt{TS1}, \texttt{OT2}.
    	\item Слово \textsf{б} значајно побољшано у свим стиловима.
    	\item Додата експериментална кодирања \texttt{ECMSRB1} и \texttt{ECMSRB2}.
    \end{itemize}
    
    \subsection*{3.1}
    
    \begin{itemize}
    	\item Исправљени багови.
    	\item Опције \textsf{noint} и \textsf{nosymb} су додате.
    	\item Измењени знаци $\leq$, $\geq$, $\nleq$, $\ngeq$.
    \end{itemize}
    
    \subsection*{3.0}
    
    \begin{itemize}
    	\item Слово б је измењено.
    	\item Додата конверзија из ћирилице у латиницу.
    	\item Математички фонтови су измењени (\textsf{nomath} опција је активна).
    	\item Додат је усправни интеграл.
    \end{itemize}
    
    \subsection*{2.0}
    
    \begin{itemize}
    	\item Додата подршка за слова {\fontencoding{OT2}\selectfont \'g, \'k, \'s, \'z}.
    \end{itemize}
    
    \subsection*{1.1}
    
    \begin{itemize}
    	\item Неколико \verb|.map| фајлова спојено у један.
    \end{itemize}

	\pagebreak

	\appendix
	
	\section{Кодирање \textit{ECMSRB1}}\label{ecs1}
	
	\fonttable{ecmsrb1cmsrbrr}
	
	\pagebreak
	
	\section{Кодирање \textit{ECMSRB2}}\label{ecs2}
	
	\fonttable{ecmsrb2cmsrbrr}
	
\end{document}
